% ****** Start of file apssamp.tex ******
%
%   This file is part of the APS files in the REVTeX 4.2 distribution.
%   Version 4.2a of REVTeX, December 2014
%
%   Copyright (c) 2014 The American Physical Society.
%
%   See the REVTeX 4 README file for restrictions and more information.
%
% TeX'ing this file requires that you have AMS-LaTeX 2.0 installed
% as well as the rest of the prerequisites for REVTeX 4.2
%
% See the REVTeX 4 README file
% It also requires running BibTeX. The commands are as follows:
%
%  1)  latex apssamp.tex
%  2)  bibtex apssamp
%  3)  latex apssamp.tex
%  4)  latex apssamp.tex
%
\documentclass[%
 reprint,
%superscriptaddress,
%groupedaddress,
%unsortedaddress,
%runinaddress,
%frontmatterverbose, 
%preprint,
%preprintnumbers,
%nofootinbib,
%nobibnotes,
%bibnotes,
 amsmath,amssymb,
 aps,
%pra,
%prb,
%rmp,
%prstab,
%prstper,
%floatfix,
]{revtex4-2}

\usepackage{graphicx}% Include figure files
\usepackage{dcolumn}% Align table columns on decimal point
\usepackage{bm}% bold math
%\usepackage{hyperref}% add hypertext capabilities
%\usepackage[mathlines]{lineno}% Enable numbering of text and display math
%\linenumbers\relax % Commence numbering lines

%\usepackage[showframe,%Uncomment any one of the following lines to test 
%%scale=0.7, marginratio={1:1, 2:3}, ignoreall,% default settings
%%text={7in,10in},centering,
%%margin=1.5in,
%%total={6.5in,8.75in}, top=1.2in, left=0.9in, includefoot,
%%height=10in,a5paper,hmargin={3cm,0.8in},
%]{geometry}

\begin{document}

\preprint{APS/123-QED}

\title{Exact Thermalization in Lorentzian Environemts}% Force line breaks with \\
\thanks{A footnote to the article title}%

\author{Ludovico Lenzuni}
 %\altaffiliation[Also at ]{Physics Department, XYZ University.}%Lines break automatically or can be forced with \\
\author{Vasco Cavina}%
 %\email{Second.Author@institution.edu}
\affiliation{Scuola Normale Superiore}%

%\collaboration{MUSO Collaboration}%\noaffiliation

\author{Gonzalo Manzano}
 \homepage{http://www.Second.institution.edu/~Charlie.Author}
\affiliation{
 IFISC}%
%\affiliation{
 %Third institution, the second for Charlie Author
%}%
%\author{Delta Author}
\%affiliation{%
% Authors' institution and/or address\\
 This line break forced with \textbackslash\textbackslash
}%

%\collaboration{CLEO Collaboration}%\noaffiliation

\date{\today}% It is always \today, today,
             %  but any date may be explicitly specified

\begin{abstract}
In this work we study the behavior of a qubit coupled to an environment composed of a continuum of bosonic modes producing a Lorentzian spectrum with peak resonant to the qubit frequency. This will produce non-Markovian dynamics on the qubit. We show that starting the environment in a thermal state causes the qubit to exactly thermalize in a finite time $t^*$ under the hypothesis of low temperature and weak coupling. We give an estimate of how $t^*$ depends on the various parameters and characterize the critical coupling under which this mechanism stops. 
%\begin{description}
%\item[Usage]
%Secondary publications and information retrieval purposes.
%\item[Structure]
%You may use the \texttt{description} environment to structure your abstract;
%use the optional argument of the \verb+\item+ command to give the category of each item. 
%\end{description}
\end{abstract}

%\keywords{Suggested keywords}%Use showkeys class option if keyword
                              %display desired
\maketitle

%\tableofcontents

\section{INTRODUCTION}

\section{MODEL}
We consider a qubit coupled to an environment of bosonic modes. The total Hamiltonian can be written as 
\begin{eqnarray}
H &&\equiv  H_S + H_E + V \; ,
\\
H_S &&\equiv \frac{\omega_S}{2} \sigma_z \; ,
\\
H_E &&\equiv \sum_k \omega_k \,  a^\dagger _k a_k \; ,
\\
V &&\equiv \sum_k \; \Big[ g_k \, \sigma^- a^\dagger_k + g_k^* \,\sigma^+ a_k \big) \Big] \, . 
\end{eqnarray}
Here the introduced creation and annihilation operators satisfy standard commutation relations $\{ a_k^\dagger, \; a_{k'} \}= \delta_{kk'}$. \\
This type of quibit-mode coupling is usually referred as Jaynes-Cummings (JC) coupling and is equivalent to the standard coupling $V_x= \sigma_x \otimes X$ under the rotating-wave approximation (RWA). We can also define the spectral density as
\begin{equation}
    J(\omega)= \sum_k|g_k|^2 \delta (\omega-\omega_k) \; .
\end{equation}
We now assume that the spectral density of $E $ is Lorentzian while the one of $B$ is flat (white noise):
\begin{eqnarray}
    J_E(\omega)= \frac{ \gamma}{ (\omega - \omega_S) ^2 + \gamma^2} \; .
\end{eqnarray}
Since we are interested in the study of a thermalization process, we impose that the initial state of the environment is thermal with inverse temperature $\beta$. In the coninuous limit this translate to the KMS condition on the expectation values. \\
Now we introduce a mapping very similar to the standard RC mapping that will be useful to prove our result. As proven by \cite{Tamascelli2017-hx}, two bosonic environments produce the same dynamics on a system if a correspondence can be found between the operators with which they couple to the system has the same time-dependent expectation values and two-times correlations.\\
We now introduce a new model on which we aim to do the mapping with the current one. We consider a qubit coupled to an auxiliary system which is in turn weakly coupled to a residual bath of oscillators.
\begin{eqnarray}
    H &&= H_S + H_R + H_B + V_{SR} + V_{RB} \; ,
    \\
    H_R &&= \omega_0 \: a_R^\dagger a _R \; ,
    \\
    H_B &&\equiv \sum_j \omega_j \,  a^\dagger _j a_j \; ,
    \\
    V_{SR} &&=  c \, \sigma^- a^\dagger_R + c^* \sigma^+ a_R \; ,
    \\
    V_{SB} &&=  \Sigma_j \big[ h_j a_R a^\dagger_j + h_j^* a_R^\dagger a_j  \big].    
\end{eqnarray}
We will assume both $R$ and $B$ in an initial thermal state. Besides, we will suppose the coupling between $R$ and $B$ to be sufficiently small that the dynamics induced by $B$ on the residual system can be modeled by the thermal markovian GKSL generator $\mathcal{L_\beta}$:
\begin{eqnarray}   
    \frac{d}{dt} \rho_{SR}(t) &&= \mathcal{L_\beta}[\rho_{SR}(t)] \\
    &&\equiv -i [H_S+V_{SR} + H_R, \rho_{SR}(t)]+ \gamma \big( n_\beta(\omega_0) \mathcal{D}_a [\rho_{SR}(t)] +(n_\beta(\omega_0) +1)\mathcal{D}_{a^\dagger} [\rho_{SR}(t)] \big) 
\end{eqnarray}
Therefore, quantity we are interested in are:
\begin{eqnarray}
    C_R^+(t) &&= |c|^2 \langle a_R(t) a^\dagger (0) \rangle \; ,
    \\
    C_R^- (t) &&= |c|^2 \langle a_R^\dagger(t) a (0) \rangle \; .
    
\end{eqnarray}
and the ones of the original model 
\begin{eqnarray}
    C_E^+(t) &&=  \langle A_E(t) A^\dagger (0) \rangle \; ,
    \\
    C_E^- (t) &&=  \langle A_E^\dagger(t) A (0) \rangle \; .
    
\end{eqnarray}
It can be shown that $\langle a_R (t) \rangle = \langle a^\dagger_R (t) \rangle = \langle A_E (t) \rangle = \langle A^\dagger_E (t) \rangle =0$ and the same is true for homogeneous correlators $\langle a_R (t) a_R (0) \rangle = \langle a^\dagger_R (t) a^\dagger_R (0) \rangle = \langle A_E (t) A_E (0) \rangle = \langle A^\dagger_E (t) A^\dagger_E (0) \rangle =0$.\\

\bibliography{apssamp}% Produces the bibliography via BibTeX.

\end{document}
%
% ****** End of file apssamp.tex ******
